\documentclass[parskip=half-, 12pt, DIV=15, twoside=semi]{scrbook}
\usepackage[sans, frenchstyle]{kpfonts-otf}
%\usepackage{showframe}
\usepackage{polyglossia}
\setmainlanguage{french}
\usepackage{xcolor}
\usepackage{tabularray}
\usepackage{multicol}
\usepackage{enumitem}
\setlist[enumerate]{label=\textbf{\alph*)}}
\SetEnumitemKey{cols}
{
    before = \begin{multicols}{#1},
    after = \end{multicols}
}
\usepackage[use-files]{xsim}
\xsimsetup{
    path=xsim,
    load-style=mytemplate,
    exercise/template=mytemplate,
}
\begin{document}

\begin{huge}
    Chapitre 1 : les entiers
\end{huge}

\medbreak

L'exercice 1 est à faire sur cette feuille. Les suivants sont à faire sur feuille annexe. Pour chaque calcul à effectuer, recopie-le d'abord sur ta feuille. 

Utilise des couleurs, souligne les éléments importants, tout ce qui pourra t'aider à retenir et progresser.

\medbreak

\begin{exercise}
    Complète le tableau en calculant la valeur numérique des expressions suivantes.
    \begin{center}
        \(
        \begin{tblr}{
                hlines,
                vlines,
                colspec={*{3}{c}*{2}{X[c]}}, 
                rows={2.5cm,m},
                row{1}={1cm, gray!40},
                column{1-3}={gray!40}
            }
            a  & b  & c  & a+b-c & a-(b+c) \\
            10 & -3 & 8  &       &         \\
            -6 & -5 & 2  &       &         \\
            3  & -8 & -2 &       &         \\
            7  & -2 & -5 &       &         \\
        \end{tblr}
        \)
    \end{center}
\end{exercise}

\begin{exercise}
    Effectue en une seule étape les puissances suivantes.
    \begin{enumerate}[cols=6, itemsep=1em]
        \item \((-3)^2\)
        \item \(-1^6\)
        \item \(-5^2\)
        \item \((-3)^3\)
        \item \((-2)^4\)
        \item \(-1^5\)
        \item \(3^0\)
        \item \(25^1\)
        \item \((-8)^0\)
        \item \(0^7\)
        \item \(3^{3} \)
        \item \(7^{1} \)
        \item \(( - 2)^{4} \)
        \item \({- 1}^{6} \)
        \item \(( - 1)^{6} \)
        \item \(( - 2)^{3} \)
        \item \((-5)^3\)
        \item \(-3^2\)
    \end{enumerate}
\end{exercise}

\pagebreak
\begin{exercise}
    Effectue. 
    \begin{enumerate}[cols=3, itemsep=1em]
        \item $(-7)+(-10)$
        \item $( + 7) + ( - 10)$
        \item $25 - ( + 7)$
        \item $-50+(-2)-(-9)$
        \item $-2+5-(-7)$
        \item $14-(-5)+(-7)-(+10)$
        \item $12 + 7 - 25 + 3$
        \item $10 + ( - 7) - ( - 15)$
        \item $- 14 + 7$
        \item $( - 5) + ( + 8) - ( + 9)$
        \item $14 - ( - 7) + 17 - 28 - 5$
        \item $( - 5) + ( - 4) - ( - 12)$
        \item $- 5 + ( - 7) - ( + 8)$
        \item $- 12 - 13 - 14 + 15$
        \item $(-3)- (-5) - 10$
    \end{enumerate}
\end{exercise}

\begin{exercise}
    Effectue.
    \begin{enumerate}[cols=3, itemsep=.3cm]
        \item \(4\cdot (-5) \cdot 7\)
        \item \(-2\cdot (-3)\cdot (-4)\)
        \item \(2\cdot (-1)\cdot 1\)
        \item \(7\cdot (-2)\cdot (-2)\)
        \item \(-5\cdot (-2)\cdot (-2)\cdot  (-10)\)
        \item \(6\cdot (-3) \cdot 2\)
        \item \(-1 \cdot 1 \cdot (-1) \cdot (-1)\)
        \item \(3\cdot (-3) \cdot 4\)
        \item \(-2\cdot (-2) \cdot 5\)
        \item \(5 \cdot 2 \cdot (-2) \cdot (-1)\)
        \item \(-10 \cdot (-1) \cdot 7\)
        \item \(-2 \cdot 8 \cdot (-1) \cdot (-1)\)
    \end{enumerate}
\end{exercise}

\begin{exercise}
    Code/décode.
    \begin{enumerate}[cols=2]
        \item L'opposé de 7
        \item La somme de 8 et de l'opposé de 10 vaut l'opposé de 2
        \item La différence entre l'opposé de 6 et l'opposé de 7
        \item $-8+2=-6$
        \item $-5$
        \item $(-4)^2$
        \item $-4^2$
        \item $-3+3^3=24$
    \end{enumerate}
\end{exercise}

\begin{exercise}
    Place les points suivants sur une droite que tu gradueras de sorte que \emph{2cm = 1 unité}.
    \begin{enumerate}[label=\ ]
        \item $A$ si tu sais que son abscisse vaut -3
        \item $B$ si tu sais que son abscisse vaut 0,5
        \item $C$ tel que $abs C = -2.5$
        \item $D$ si tu sais que son abscisse vaut la somme des abscisses de $B$ et $C$
        \item $E$ si tu sais que son abscisse vaut 3,5
        \item $F$ si tu sais que les abscisses de $E$ et $F$ sont des nombres opposés.
    \end{enumerate}
\end{exercise}

\pagebreak
\begin{exercise}
    Effectue.
    \begin{enumerate}[cols=2, itemsep=1emm]
        \item \(3^{2} + 5 \cdot ( - 2) \)
        \item \(- 15 + ( - 3 + 7)^{2} \)
        \item \(\left( - 3 + ( - 2) \right) \cdot (20 - 12) \)
        \item \(- 40 + ( - 2) \cdot ( - 8) \)
        \item \(14 \cdot ( - 2) + ( - 20) \)
        \item \(( - 3) + ( - 8) - ( - 8) \cdot 2 \)
        \item \(( - 3)^{3} + ( - 2) \cdot 8 \)
        \item \(( - 3)^{2} + \left( ( - 3) \cdot ( - 2) + 7 \right) \)
    \end{enumerate}
\end{exercise}

\begin{exercise}
    Place les élements suivants sur une droite graduée en graduant de sorte que $1cm = 1$ unité. \\

    $\hfil abs~A = 3 \hfil abs~B=-2 \hfil abs~C=-1 \hfil abs~D=2 \hfil abs~E=5$
\end{exercise}

\begin{exercise}
    Effectue.
    \begin{enumerate}[itemsep=.4cm, cols=2]
        \item   \(3^{2} \cdot ( - 3) + 2^{3} \cdot ( - 2)^{2} \)
        \item \(- 1 \cdot ( - 30) \cdot ( - 2)^{4} \)
        \item \(12 \cdot ( - 3)^{2} \)
        \item \(6^{2} \cdot ( - 2) - 2 \cdot 4^{2} \)
        \item \(3 \cdot ( - 5) + 5^{2} \cdot ( - 1) \)
        \item \(4 \cdot ( - 7)^{2} - 1 \cdot ( - 6) \)
        \item \(2 \cdot  2^{3} \cdot ( - 2)^{3} \)
        \item \(13 \cdot 10^{2} - 2 \cdot ( - 2) \)
        \item \(9 \cdot ( - 4) + 2^{3} \cdot 1 \)
        \item \(- 3\  \cdot ( - 6)^{2} + ( - 1)^{6} \)
        \item \( ( - 4)^{2} + 7 \cdot ( - 20) \)
        \item \(3^{3} \cdot ( - 2) - 14 \cdot ( - 2) \)
        \item \(-7 + 7 \cdot 2 - (-3)^2\)
        \item \(5 \cdot (-5) \cdot (-2) + \left(30:(-5)\cdot7- (-2)\right)\)
    \end{enumerate}
\end{exercise}

\begin{exercise}
    Traduis les phrases suivantes par un calcul puis effectue-le.
    \begin{enumerate}[itemsep=.3cm]
        \item La somme de l'opposé de \(4\) et de l'opposé de \(5\).
        \item Le produit de la différence entre \(4\) et \(7\) par l'opposé de \(9\).
        \item La différence entre \(5\) et le produit de \(7\) par l'opposé de \(2\).
        \item Le produit de la somme de \(5\) et de l'opposé de \(2\) par la différence entre \(3\) et \(10.\)
        \item La somme du double de l'opposé de 4 et de 5.
    \end{enumerate}
\end{exercise}

\begin{exercise}
    Calcule la valeur numérique des expressions suivantes si tu sais que
    \(a = 2\), \(b = - 3\) et \(c = - 1\).
    \begin{enumerate}[cols=3, itemsep=.4cm]
        \item \(ab + c\)
        \item \(3b - 7a\)
        \item \((c - b) + a\)
        \item \(3bc - b^{2}\)
        \item \(a + 2b - c\)
        \item \(b \cdot ( - c)\)
        \item \(-2b+c-(10-3)\)
        \item \(a^3 - b^2 - c^6\)
    \end{enumerate}
\end{exercise}

\begin{exercise}
    Vrai ou faux~? Justifie dans tous les cas.
    \begin{enumerate}
        \item \(( - 4) \cdot a\) est toujours négatif.
        \item \(a^{2}\) est toujours positif.
        \item Le produit de \(a\) par son opposé est négatif.
        \item Le double de \(a\) est positif.
        \item \(\dfrac{-3}{7}\) est l'inverse de \(frac{3}{7}\)
    \end{enumerate}
\end{exercise}

\end{document}